Fix \(\mathcal M\in\mathfrak M_{+}\), \(x\in\mathcal X\), and \(y\in\mathcal Y\). By strict primitive positivity, for each \(e\in\mathcal E\),
\[
\rho_x(e)=\sum_{u\in\mathcal U}a_x(u)S_{ue}>0,
\tag{1}
\]
so every conditional probability below is defined. Summing the source factorization in Assumption~\ref{ass:latent-shift-factorization} over the outcome and then conditioning on \((E=e,X=x)\) gives
\[
\begin{aligned}
B_x(w,e)
&=\frac{\sum_{u\in\mathcal U}S_{ue}a_x(u)M_{wu}}{\rho_x(e)}\\
&=\sum_{u\in\mathcal U}M_{wu}R_x(u,e).
\end{aligned}
\tag{2}
\]
Thus \(B_x=MR_x\). By target-mechanism invariance and the definition of \(b\),
\[
b_w=\sum_{u\in\mathcal U}M_{wu}q_u,
\tag{3}
\]
whence \(b=Mq\).

Retaining the outcome factor in the same source calculation yields
\[
\begin{aligned}
c_{x,y}(e)
&=\frac{\sum_{u\in\mathcal U}\sum_{w\in\mathcal W}
 S_{ue}a_x(u)M_{wu}f_{x,y}(u,w)}{\rho_x(e)}\\
&=\sum_{u\in\mathcal U}R_x(u,e)
   \sum_{w\in\mathcal W}M_{wu}f_{x,y}(u,w)\\
&=\sum_{u\in\mathcal U}g_{x,y}(u)R_x(u,e).
\end{aligned}
\tag{4}
\]
Therefore \(c_{x,y}=g_{x,y}^{\top}R_x\). Under \(\operatorname{do}(X=x)\), the graph-truncated target factorization deletes the treatment kernel and retains the invariant \(U\), \(W\), and \(Y\) mechanisms. Consequently,
\[
\begin{aligned}
\theta_{x,y}
&=\sum_{u\in\mathcal U}\sum_{w\in\mathcal W}
 q_uM_{wu}f_{x,y}(u,w)\\
&=\sum_{u\in\mathcal U}q_ug_{x,y}(u)
=g_{x,y}^{\top}q.
\end{aligned}
\tag{5}
\]
This uses Assumption~\ref{ass:target-mechanism-invariance}; in particular, (5) is a causal intervention identity rather than a definition in terms of the observed law.

Now take \(\lambda\in\Lambda_x(\mathcal D)\). By Definition~\ref{def:balancing-fiber}, \(B_x\lambda=b\). Equations (2)--(3) imply
\[
M(R_x\lambda-q)=0.
\tag{6}
\]
Membership in \(\mathfrak M_{+}\), through Definition~\ref{def:positive-latent-shift-class}, includes \(\operatorname{rank}(M)=r\). Hence the linear map \(z\mapsto Mz\) on \(\mathbb R^r\) is injective, and (6) gives
\[
R_x\lambda=q.
\tag{7}
\]
Put \(p_x(e)=P_O(E=e,X=x)=\pi_e\rho_x(e)\). Strict positivity and (1) give \(p_x(e)>0\), so Definition~\ref{def:unconditional-weight-map} is well defined and gives \(\kappa(e)=\lambda(e)/p_x(e)\). By Definition~\ref{def:unconditional-balancing-moments},
\[
H_x(w,e)=p_x(e)B_x(w,e),\qquad
z_{x,y}(e)=p_x(e)c_{x,y}(e).
\tag{8}
\]
It follows coordinatewise that
\[
(H_x\kappa)_w
=\sum_{e\in\mathcal E}p_x(e)B_x(w,e)\frac{\lambda(e)}{p_x(e)}
=(B_x\lambda)_w=b_w,
\tag{9}
\]
and, using (4), (5), and (7),
\[
z_{x,y}^{\top}\kappa
=\sum_{e\in\mathcal E}c_{x,y}(e)\lambda(e)
=g_{x,y}^{\top}R_x\lambda
=g_{x,y}^{\top}q
=\theta_{x,y}.
\tag{10}
\]
Equations (2)--(5), (9), and (10) are all the asserted identities.